\scanpage{6}
\diagram{\begin{tikzcd}[ampersand replacement=\&,row sep=small,column sep=small]
\& \pi \arrow[d,"u"'] \\
\& G(\unsure{Q|\mathbb Q}_p) \arrow[d,phantom,"\rotatebox{270}{$\subset$}" description] \\
\& G(C) \arrow[d,phantom,"\unsure{\wr\vert|\simeq|j_1}" description] \\
1 \arrow[r,dash] \& G'(C) \arrow[r] \arrow[d,phantom,"\rotatebox{90}{$\subset$}" description] \& E \arrow[r] \& \pi \arrow[l,bend right=32,"\sigma_1"'] \arrow[l,bend left=32,"\sigma_2"] \arrow[r,dash] \& 1 \\
\& G'(Q_p) \\
\& G_m(Q_p)=Q_p^* \arrow[u,hook,"j_2"] \\
\& \pi \arrow[u,"\chi"]
\end{tikzcd}}

\deleted{\unsure{$\alpha$}} $\alpha$ permet d'identifier une certaine ext. de $E$ de $\pi$ par $G'(C)$ avec une \deleted{autre} ext. de $\pi$ par $G(C)$

Le splitting $\sigma_2$ de la première est tel que $\sigma_2(\pi)$ commute à $G'(Q_p)$ et a fortiori à $j_2\chi(\pi)$. On trouve donc
\[
\left\{\begin{aligned}
\gamma&\longmapsto\widetilde{\sigma}_2(\gamma)=\sigma_2(\gamma).j_2(\chi(\gamma))\\
\pi&\longrightarrow E
\end{aligned}\right.
\]
un autre splitting.

Le splitting $\sigma_1$ \deleted{\unsure{de|dans}} la seconde commute à $G(Q_p)$, et à fortiori à $u(\pi)$. On \deleted{trouve donc}
\[
\gamma\longmapsto\sigma_1(\gamma)u(\gamma)=\widetilde{\sigma}_1(\gamma)
\]

On écrit que $\widetilde{\sigma}_1(\gamma)=\widetilde{\sigma}_2(\gamma)$

\margin{\unsure{$\lrcorner$|1}}% margin note written sideways
i.e.
\[
\sigma_2(\gamma)\,j_2(\chi(\gamma))=\sigma_1(\gamma)\,u(\gamma)
\]
donc
\[
\underbrace{\sigma_1(\gamma)^{-1}\sigma_2(\gamma)}_{\varphi(\gamma)}
=\deleted{\unsure{j_2}}\underset{\substack{\rotatebox{270}{$\in$}\\G(Q_p)}}{u(\gamma)}\mathord{.}
\underset{\substack{\rotatebox{270}{$\in$}\\j_2(Q_p^*)}}{\underline{j_2(\chi(\gamma^{-1}))}}
\]

1\textsuperscript{o}) Sur $\operatorname{Ker}\chi$, on a $\varphi(\gamma)=u(\gamma)$. donc $u$ connu sur $\operatorname{Ker}\chi$

2\textsuperscript{o}) $\operatorname{Ker}u=\deleted{\chi(\gamma)}\{\gamma,\chi(\gamma)=1,\varphi(\gamma)=1\}$ connu. Soit $\Pi_1=\pi/\operatorname{Ker}u$. C'est un gpe analytique $p$-adique, et $\Pi_1\xrightarrow{\chi_1}Q_p^{\unsure{*|1}}$ induit par $p$
% archive sheet number: 5
